\subsection{Diagnosing a completed fit}
\label{sec:pfit-diagnose}

After a fitting run has finished, the user requests a post-fit assessment with

\begin{verbatim}
/pfit-diagnose <session-name> <run-id>
\end{verbatim}

The diagnosis step evaluates a completed run. It does not refit the model and
it does not inspect mutable working files in place of the run that was actually
performed. The agent resolves one run directory, reads the configuration and
model snapshots saved with that run, and bases the assessment on the artifacts
written by the optimizer.

The purpose of this step is to decide whether the fit is usable, and if not,
what should be changed before rerunning. The agent examines the global-search
and refinement logs, the final parameter values, the sloppiness report, the
saved trajectories, residuals, and the measured-versus-fitted plots. The plots
are part of the evidence: a small numerical objective is not sufficient if the
curves reveal a wrong observable mapping, a missed transient, or systematic
structure in the residuals.

The report is written to

\begin{verbatim}
sessions/<session-name>/outputs/<run-id>/fit_diagnosis.txt
\end{verbatim}

and gives a single supported verdict where the evidence allows one. Typical
findings include failed integrations, insufficient global-search budget,
premature population collapse, a gradient stage that stopped before convergence,
parameters pinned at bounds, poorly scaled residuals, or practically
non-identifiable parameter combinations. Each finding is tied to the run
evidence that supports it, such as loss histories, fitted values relative to
bounds, residual patterns, gradient norms, or the sloppiness spectrum.
When the refinement stage reports a JAX, Diffrax, Equinox, Lineax, NaN, or
non-finite-gradient failure, the agent should also run
\texttt{tools/autodiff\_diagnose.py}. That probe checks whether the completed
fit can be differentiated locally, compares autodiff gradients with finite
differences, tests nearby perturbations for solver walls, and scans the
generated model for non-smooth event-like logic.

Recommendations are proposed as edits to the working session, not to the saved
run snapshot. The user chooses which recommendations to apply. Depending on the
finding, the next run may be a gradient-only refinement from the diagnosed design
point or a full refit. Changes to the equations, bounds, log scaling, solver
family, or population-search settings generally require the checked translation
and full fitting workflow to be repeated.

The user should treat \texttt{/pfit-diagnose} as the point at which numerical
results are converted into scientific evidence. Its output should be reviewed
before reporting fitted parameters, interpreting identifiability, or deciding
that additional experiments are needed.
